
% ============================================================
% exercise-tags-and-stubs-v2.tex
% Include file for the Learning Abstract Mathematics project.
%
% PURPOSE
%   - Defines searchable technique/dependency tag helpers.
%   - Defines sourced exercise-record stubs:
%       * exerciseconstruction      (Types 4, 5)
%       * exerciseproofrecord       (Types 1, 2, 3, 6)
%   - Defines generated exercise-record stubs:
%       * generatedexerciseconstruction
%       * generatedexerciseproof
%   - Defines a quick generated exercise macro:
%       * \exstubgenfull
%
% DESIGN GOALS
%   - No \usepackage calls here.
%   - No theorem-environment redefinitions.
%   - Safe to \input into main.tex after your existing preamble.
%   - Uses only packages already present in your main.tex:
%       xparse, etoolbox, expl3, tcolorbox, hyperref
%
% SEARCHABILITY
%   Write literal tags in source, e.g.
%       tech:triangle-inequality
%       tech:epsilon-splitting
%       def:metric
%       thm:triangle-inequality
%   so that grep/search works on the .tex source directly.
%
% OPTIONAL INDEXING
%   If you later load makeidx/imakeidx and define \index,
%   \TechniqueEntry{...} and \DependencyEntry{...} will add index entries.
%   Otherwise they simply typeset the tag.
% ============================================================

\makeatletter

% ------------------------------------------------------------
% Optional index hooks
% ------------------------------------------------------------
\providecommand{\ExerciseTechniqueIndexName}{tech}
\providecommand{\ExerciseDependencyIndexName}{dep}

\newcommand{\exercise@maybeindextech}[1]{%
  \@ifundefined{index}{}{%
    \index[tech]{#1}%
  }%
}
\newcommand{\exercise@maybeindexdep}[1]{%
  \@ifundefined{index}{}{%
    \index[dep]{#1}%
  }%
}

% ------------------------------------------------------------
% Tag entry macros
%   IMPORTANT: pass the FULL literal tag string.
%   Example:
%     \TechniqueEntry{tech:triangle-inequality}
%     \DependencyEntry{def:metric}
% ------------------------------------------------------------
\providecommand{\TechniqueEntry}[1]{%
  \texttt{#1}%
  \exercise@maybeindextech{#1}%
}

\providecommand{\DependencyEntry}[1]{%
  \texttt{#1}%
  \exercise@maybeindexdep{#1}%
}

% ------------------------------------------------------------
% Convenience wrappers for common tag namespaces
% ------------------------------------------------------------
\providecommand{\tech}[1]{\TechniqueEntry{tech:#1}}
\providecommand{\dep}[1]{\DependencyEntry{#1}}

% ------------------------------------------------------------
% Metadata storage (sourced exercises)
% ------------------------------------------------------------
\providecommand{\CurrentExerciseAuthor}{}
\providecommand{\CurrentExerciseBook}{}
\providecommand{\CurrentExerciseID}{}
\providecommand{\CurrentExerciseType}{}
\providecommand{\CurrentExerciseTechniques}{}
\providecommand{\CurrentExerciseDependencies}{}

% ------------------------------------------------------------
% Metadata storage (generated exercises)
% ------------------------------------------------------------
\providecommand{\CurrentGeneratedExerciseID}{}
\providecommand{\CurrentGeneratedExerciseTopic}{}
\providecommand{\CurrentGeneratedExerciseType}{}
\providecommand{\CurrentGeneratedExerciseTechniques}{}
\providecommand{\CurrentGeneratedExerciseDependencies}{}

% ------------------------------------------------------------
% Per-author counters
% ------------------------------------------------------------
\newcounter{exerciseAbbott}
\newcounter{exerciseTao}
\newcounter{exerciseLebl}
\newcounter{exerciseBruckner}
\newcounter{exerciseBartleSherbert}
\newcounter{exerciseGeneric}

\newcommand{\exercise@stepauthorcounter}[1]{%
  \ifdefstrequal{#1}{Abbott}{\stepcounter{exerciseAbbott}}{%
  \ifdefstrequal{#1}{Tao}{\stepcounter{exerciseTao}}{%
  \ifdefstrequal{#1}{Lebl}{\stepcounter{exerciseLebl}}{%
  \ifdefstrequal{#1}{Bruckner}{\stepcounter{exerciseBruckner}}{%
  \ifdefstrequal{#1}{Bartle \& Sherbert}{\stepcounter{exerciseBartleSherbert}}{%
  \stepcounter{exerciseGeneric}}}}}}%
}

\newcommand{\exercise@authorcountvalue}[1]{%
  \ifdefstrequal{#1}{Abbott}{\arabic{exerciseAbbott}}{%
  \ifdefstrequal{#1}{Tao}{\arabic{exerciseTao}}{%
  \ifdefstrequal{#1}{Lebl}{\arabic{exerciseLebl}}{%
  \ifdefstrequal{#1}{Bruckner}{\arabic{exerciseBruckner}}{%
  \ifdefstrequal{#1}{Bartle \& Sherbert}{\arabic{exerciseBartleSherbert}}{%
  \arabic{exerciseGeneric}}}}}}%
}

% ------------------------------------------------------------
% Metadata setter (sourced exercises)
% ------------------------------------------------------------
\NewDocumentCommand{\ExerciseMeta}{m m m m m m}{%
  % #1 author
  % #2 book
  % #3 exercise id
  % #4 type
  % #5 techniques   (literal source-searchable tags)
  % #6 dependencies (literal source-searchable tags)
  \renewcommand{\CurrentExerciseAuthor}{#1}%
  \renewcommand{\CurrentExerciseBook}{#2}%
  \renewcommand{\CurrentExerciseID}{#3}%
  \renewcommand{\CurrentExerciseType}{#4}%
  \renewcommand{\CurrentExerciseTechniques}{#5}%
  \renewcommand{\CurrentExerciseDependencies}{#6}%
  \exercise@stepauthorcounter{#1}%
}

% ------------------------------------------------------------
% Metadata setter (generated exercises)
% ------------------------------------------------------------
\NewDocumentCommand{\GeneratedExerciseMeta}{m m m m m}{%
  % #1 exercise id
  % #2 topic
  % #3 type
  % #4 techniques
  % #5 dependencies
  \renewcommand{\CurrentGeneratedExerciseID}{#1}%
  \renewcommand{\CurrentGeneratedExerciseTopic}{#2}%
  \renewcommand{\CurrentGeneratedExerciseType}{#3}%
  \renewcommand{\CurrentGeneratedExerciseTechniques}{#4}%
  \renewcommand{\CurrentGeneratedExerciseDependencies}{#5}%
}

% ------------------------------------------------------------
% Shared header (sourced exercises)
% ------------------------------------------------------------
\newcommand{\ExerciseRecordHeader}{%
\begin{tcolorbox}[colback=gray!6,colframe=gray!40,arc=2pt,
  left=6pt,right=6pt,top=4pt,bottom=4pt,
  title={\small\textbf{Exercise Record}}]
\small
\textbf{Author index number. }
\CurrentExerciseAuthor~\exercise@authorcountvalue{\CurrentExerciseAuthor}

\medskip
\textbf{Source. }
\CurrentExerciseAuthor, \textit{\CurrentExerciseBook}, \CurrentExerciseID

\medskip
\textbf{Type. }
\CurrentExerciseType

\medskip
\textbf{Technique tags. }
\texttt{\CurrentExerciseTechniques}

\medskip
\textbf{Dependency tags. }
\texttt{\CurrentExerciseDependencies}
\end{tcolorbox}
}

% ------------------------------------------------------------
% Shared header (generated exercises)
% ------------------------------------------------------------
\newcommand{\GeneratedExerciseRecordHeader}{%
\begin{tcolorbox}[colback=gray!6,colframe=gray!40,arc=2pt,
  left=6pt,right=6pt,top=4pt,bottom=4pt,
  title={\small\textbf{Exercise Record}}]
\small
\textbf{Origin. } Generated exercise

\medskip
\textbf{Exercise ID. } \CurrentGeneratedExerciseID

\medskip
\textbf{Topic. } \CurrentGeneratedExerciseTopic

\medskip
\textbf{Type. } \CurrentGeneratedExerciseType

\medskip
\textbf{Technique tags. } \texttt{\CurrentGeneratedExerciseTechniques}

\medskip
\textbf{Dependency tags. } \texttt{\CurrentGeneratedExerciseDependencies}
\end{tcolorbox}
}

% ------------------------------------------------------------
% Stub A — Construction / Find / Compute (sourced)
% ------------------------------------------------------------
\NewDocumentEnvironment{exerciseconstruction}{}{%
  \subsection*{Exercise}
  \ExerciseRecordHeader

  % Suggested searchable metadata block:
  % tech:...
  % def:...
  % thm:...

  \begin{remark*}[Return]
  $\leftarrow$ Back-link to the relevant notes section.
  \end{remark*}

  \begin{remark*}[Task]
  State the exercise verbatim here.
  \end{remark*}

  \begin{remark*}[Construction]
  State the object or computed value explicitly before any justification.
  \end{remark*}

  \begin{proof}
}{%
  \end{proof}

  \begin{remark*}[Dependencies]
  Expand the dependency tags into prose if helpful.
  \end{remark*}

  \begin{remark*}[Proof shape]
  Record the structural pattern used:
  direct / contradiction / verification / construction / computation.
  \end{remark*}

  \begin{remark*}[Reflection]
  Optional: explain why the construction works and what feature is essential.
  \end{remark*}
}

% ------------------------------------------------------------
% Stub B — Proof / Verify / Prove-or-Counterexample / Characterization (sourced)
% ------------------------------------------------------------
\NewDocumentEnvironment{exerciseproofrecord}{}{%
  \subsection*{Exercise}
  \ExerciseRecordHeader

  % Suggested searchable metadata block:
  % tech:...
  % def:...
  % thm:...

  \begin{remark*}[Return]
  $\leftarrow$ Back-link to the relevant notes section.
  \end{remark*}

  \begin{remark*}[Task]
  State the exercise verbatim here.
  \end{remark*}

  \begin{remark*}[Strategy]
  Briefly describe the proof plan.
  \end{remark*}

  \begin{proof}
}{%
  \end{proof}

  \begin{remark*}[Dependencies]
  Expand the dependency tags into prose if helpful.
  \end{remark*}

  \begin{remark*}[Proof shape]
  Record the structural pattern:
  direct / contradiction / case split / subsequence extraction / squeeze / iff.
  \end{remark*}
}

% ------------------------------------------------------------
% Stub A — Construction / Find / Compute (generated)
% ------------------------------------------------------------
\NewDocumentEnvironment{generatedexerciseconstruction}{}{%
  \subsection*{Exercise}
  \GeneratedExerciseRecordHeader

  % Suggested searchable metadata block:
  % tech:...
  % def:...
  % thm:...

  \begin{remark*}[Task]
  State the generated exercise here.
  \end{remark*}

  \begin{remark*}[Construction]
  State the object or computed value explicitly before justification.
  \end{remark*}

  \begin{proof}
}{%
  \end{proof}

  \begin{remark*}[Dependencies]
  Expand the dependency tags into prose if helpful.
  \end{remark*}

  \begin{remark*}[Proof shape]
  Record the structural pattern used:
  direct / contradiction / verification / construction / computation.
  \end{remark*}

  \begin{remark*}[Reflection]
  Optional: explain why the construction works.
  \end{remark*}
}

% ------------------------------------------------------------
% Stub B — Proof / Verify / Prove-or-Counterexample / Characterization (generated)
% ------------------------------------------------------------
\NewDocumentEnvironment{generatedexerciseproof}{}{%
  \subsection*{Exercise}
  \GeneratedExerciseRecordHeader

  % Suggested searchable metadata block:
  % tech:...
  % def:...
  % thm:...

  \begin{remark*}[Task]
  State the generated exercise here.
  \end{remark*}

  \begin{remark*}[Strategy]
  Briefly describe the proof plan.
  \end{remark*}

  \begin{proof}
}{%
  \end{proof}

  \begin{remark*}[Dependencies]
  Expand the dependency tags into prose if helpful.
  \end{remark*}

  \begin{remark*}[Proof shape]
  Record the structural pattern:
  direct / contradiction / case split / squeeze / iff.
  \end{remark*}
}

% ------------------------------------------------------------
% Quick unsolved generated exercise stub
% ------------------------------------------------------------
\newcommand{\exstubgenfull}[5]{%
  \GeneratedExerciseMeta{#1}{#2}{GENERATED}{#3}{#4}%

  \subsection*{Exercise}
  \GeneratedExerciseRecordHeader

  \begin{remark*}[Task]
  #5
  \end{remark*}

  \begin{remark*}[Work Area]
  Write your proof or computation here.
  \end{remark*}

  \bigskip
}

% ------------------------------------------------------------
% Controlled vocabularies
% ------------------------------------------------------------

% TECHNIQUE TAGS (recommended)
% tech:unpack-definition
% tech:triangle-inequality
% tech:epsilon-splitting
% tech:tail-argument
% tech:choose-N
% tech:case-splitting
% tech:contradiction
% tech:subsequence-extraction
% tech:diagonal-argument
% tech:squeeze
% tech:monotone-bound
% tech:supremum-characterization
% tech:metric-verification
% tech:ball-containment
% tech:diameter-bound
% tech:compactness-argument

% DEPENDENCY TAGS (recommended)
% def:metric
% def:open-ball
% def:bounded
% def:diameter
% def:convergence
% def:cauchy-sequence
% def:supremum
% def:infimum
% thm:triangle-inequality
% thm:archimedean-property
% thm:monotone-convergence
% thm:bolzano-weierstrass
% thm:squeeze-theorem

% ------------------------------------------------------------
% Example usage (sourced)
% ------------------------------------------------------------
%
% % =========================================================
% % Exercise Metadata
% % tech:triangle-inequality
% % tech:epsilon-splitting
% % def:convergence
% % thm:triangle-inequality
% % =========================================================
%
% \ExerciseMeta
%   {Abbott}
%   {Understanding Analysis}
%   {Exercise 2.3.8}
%   {PROOF}
%   {tech:triangle-inequality, tech:epsilon-splitting}
%   {def:convergence, thm:triangle-inequality}
%
% \begin{exerciseproofrecord}
%
% \begin{remark*}[Technique]
% \TechniqueEntry{tech:triangle-inequality},
% \TechniqueEntry{tech:epsilon-splitting}
% \end{remark*}
%
% \begin{remark*}[Dependencies]
% \DependencyEntry{def:convergence},
% \DependencyEntry{thm:triangle-inequality}
% \end{remark*}
%
% Let $\varepsilon > 0$. Since $x_n \to x$, there exists $N$ such that
% $d(x_n,x)<\varepsilon/2$ for all $n\ge N$. If $m,n\ge N$, then
% \[
% d(x_m,x_n)\le d(x_m,x)+d(x,x_n)<\varepsilon.
% \]
%
% \end{exerciseproofrecord}
%
% ------------------------------------------------------------
% Example usage (generated)
% ------------------------------------------------------------
%
% % =========================================================
% % Exercise Metadata
% % tech:supremum-characterization
% % tech:unpack-definition
% % def:supremum
% % def:infimum
% % =========================================================
%
% \exstubgenfull
%   {EX-BOUNDS-CHATGPT-001}
%   {Bounds / Suprema / Infima}
%   {tech:supremum-characterization, tech:unpack-definition}
%   {def:supremum, def:infimum}
%   {Let $A=\{1/n : n\in\mathbb N\}$. Determine $\sup A$ and $\inf A$.}
%
% ------------------------------------------------------------

\makeatother
